% Part: first-order-logic
% Chapter: natural-deduction
% Section: quantifier-rules

\documentclass[../../../include/open-logic-section]{subfiles}

\begin{document}

\olfileid[hi]{fol}{ntd}{qrl}

\olsection{परिमाणक-नियम}

\subsection{$\lforall$ के नियम}

\begin{defish}
\AxiomC{$!A(a)$}
\RightLabel{\Intro{\lforall}}
\UnaryInfC{$\lforall[x][\Atom{!A}{x}]$}
\DisplayProof
\hfill
\AxiomC{$\lforall[x][\Atom{!A}{x}]$}
\RightLabel{\Elim{\lforall}}
\UnaryInfC{$!A(t)$}
\DisplayProof
\end{defish}

~$\lforall$ के नियमों में $t$ बंद पद है (ऐसा पद जिसमें कोई चर नहीं होता),
और $a$~ऐसा !!a{constant} है जो निष्कर्ष~$\lforall[x][!A(x)]$ में या ऐसी किसी
!!{undischarged} मान्यता में नहीं आता जो पूर्वाधार~$!A(a)$ पर समाप्त होने वाली
!!{derivation} में हो। हम $a$ को \Intro{\lforall} अनुमान का
\emph{अभिलाक्षणिक चर} कहते हैं।\footnote{ऊपर के नियम में $a$ अचर है, फिर भी
ऐतिहासिक कारणों से ``अभिलाक्षणिक चर'' शब्द प्रयुक्त होता है।}

\subsection{$\lexists$ के नियम}

\begin{defish}
\AxiomC{$\Atom{!A}{t}$}
\RightLabel{\Intro{\lexists}}
\UnaryInfC{$\lexists[x][\Atom{!A}{x}]$}
\DisplayProof
\hfill
\AxiomC{$\lexists[x][\Atom{!A}{x}]$}
\AxiomC{[$\Atom{!A}{a}$]$^n$}
\DeduceC{$!C$}
\DischargeRule{\Elim{\lexists}}{n}
\BinaryInfC{$!C$}
\DisplayProof
\end{defish}

यहाँ भी $t$ बंद पद है, और $a$ ऐसा !!a{constant} है जो पूर्वाधार
$\lexists[x][!A(x)]$ में, निष्कर्ष~$!C$ में, या (मान्यताओं $!A(a)$ को छोड़कर)
ऐसी किसी !!{undischarged} मान्यता में नहीं आता जो दोनों पूर्वाधारों पर समाप्त
होने वाली !!{derivation}s में हो। हम $a$ को \Elim{\lexists} अनुमान का
\emph{अभिलाक्षणिक चर} कहते हैं।

अभिलाक्षणिक चर न तो पूर्वाधारों में आए और न ऐसी किसी !!{undischarged} मान्यता
में जो \Intro{\lforall} या \Elim{\lexists} अनुमान के पूर्वाधारों तक पहुँचने
वाली !!{derivation}s में हो—इस शर्त को \emph{अभिलाक्षणिक-चर शर्त} कहते हैं।

\begin{explain}
स्मरण करें: जब $!A$ ऐसा !!a{formula} हो जिसमें !!{variable}~$x$ मुक्त है, तो
इसे~$!A(x)$ लिखकर दिखाते हैं। उसी संदर्भ में $!A(t)$,
संक्षेप में~$\Subst{!A}{t}{x}$ है। अतः $\Intro\lexists$ नियम को ऐसे भी लिख
सकते हैं:
\begin{prooftree}
\AxiomC{$\Subst{!A}{t}{x}$}
\RightLabel{\Intro{\lexists}}
\UnaryInfC{$\lexists[x][!A]$}
\end{prooftree}
ध्यान दें कि $t$,~$!A$ में पहले से आ सकता है; जैसे $!A$ का रूप
~$\Atom{\Obj P}{t,x}$ हो सकता है। अतः $\lexists[x][\Atom{\Obj
P}{t,x}]$ को~$\Atom{\Obj P}{t,t}$ से निष्कर्षित करना
~$\Intro\lexists$ का सही प्रयोग है---पूर्वाधार में~$t$ की एक या अधिक, पर
आवश्यक नहीं कि सभी, आवृत्तियों को आबद्ध !!{variable}~$x$ से ``बदल'' सकते हैं।
परंतु $\Intro\lforall$ और~$\Elim\lexists$ की अभिलाक्षणिक-चर शर्तें माँगती हैं
कि !!{constant}~$a$,~$!A$ में न आए। इसलिए
$\lforall[x][\Atom{\Obj P}{a,x}]$ को $\Atom{\Obj P}{a,a}$ से
~$\Intro\lforall$ द्वारा सही ढंग से निष्कर्षित नहीं कर सकते।
\end{explain}

\begin{explain}
\Intro{\lexists} और \Elim{\lforall} में कोई प्रतिबंध नहीं है और पद~$t$ कुछ भी
हो सकता है, इसलिए किसी शर्त की चिंता नहीं करनी पड़ती। दूसरी ओर,
\Elim{\lexists} और \Intro{\lforall} नियमों में अभिलाक्षणिक-चर शर्त माँगती है
कि !!{constant}~$a$ निष्कर्ष में या किसी !!{undischarged} मान्यता में कहीं न
आए। तंत्र की सुदृढ़ता सुनिश्चित करने के लिए यह शर्त आवश्यक है; अर्थात् वह
केवल !!{derive} करे: उन !!{sentence}s को, जो उन !!{undischarged} मान्यताओं से
निकलते हैं। इस शर्त के बिना निम्नलिखित अनुमत होता:
\begin{prooftree}
  \AxiomC{$\lexists[x][!A(x)]$}
  \AxiomC{$\Discharge{!A(a)}{1}$}
  \RightLabel{*\Intro{\lforall}}
  \UnaryInfC{$\lforall[x][!A(x)]$}
  \RightLabel{\Elim{\lexists}}
  \BinaryInfC{$\lforall[x][!A(x)]$}
\end{prooftree}
परंतु $\lexists[x][!A(x)] \Entails/ \lforall[x][!A(x)]$।

चूँकि परिमाणकों के निरसन-नियम !!{variable}s के स्थान पर केवल बंद पद
प्रतिस्थापित करने देते हैं, अतः प्रत्येक !!{formula}, जिसे !!{sentence}s के
किसी समुच्चय से व्युत्पन्न किया जा सकता है, स्वयं !!a{sentence} है।
\end{explain}


\end{document}
